\section{Data Collection}
\label{sec:data}

No existing dataset records how Terraform blocks change together, so we built one. \textbf{TerraDS}
\cite{buehler2025terrads}, a collection of over 62,000 open-source Terraform repositories with
metadata about their modules and resources (\textit{snapshot of 2024}), served only as the starting point
for finding candidate repositories; the selection criteria, the history mining, the queries and the
labels are our own.

\paragraph{Repository selection.} Table~\ref{tab:funnel} shows the selection funnel. We kept
repositories managing at least 20 resources of one cloud provider, AWS, the most common one in
TerraDS: a single provider gives all repositories a shared vocabulary of resource types, so rules
learned on one repository transfer to another, and smaller repositories would produce candidate
lists so short that even random ranking scores well. From the 150 most-starred of them, we kept
those with at least 50 commits touching \texttt{.tf} files and removed forks and copies (repositories
sharing their first commit), which would otherwise leak history between training and test
repositories. After mining, a repository needed at least 50 usable commits for a stable
per-repository score. 8 repositories in which more than 40\% of the usable commits only add new
blocks (tutorial collections, generated examples) were set aside: they measure what is written together rather than
change impact, and are only used as a held-out group, never for training.

\begin{table}[ht]
\begin{tabular}{@{}lr@{}}
\toprule
Step & Repositories \\
\midrule
TerraDS repositories with $\geq 20$ managed AWS resources & 9,619 \\
150 most-starred, $\geq 50$ commits touching \texttt{.tf} files & 79 \\
after removing forks and copies; mined & 78 \\
$\geq 50$ usable commits, $\leq 40\%$ add-only commits (\textbf{core}) & \textbf{53} \\
\quad set aside only for their share of add-only commits (held out) & 8 \\
\bottomrule
\end{tabular}
\centering
\small
\caption{Repository selection funnel.}
\label{tab:funnel}
\end{table}

\paragraph{History mining.} Each repository is cloned and walked once, oldest commit first, over
its last 2,000 commits on the main branch (only 5 of the 53 repositories are longer). Merge commits
are diffed against their first parent, so a merged pull request counts as one change. Every version
of every \texttt{.tf} file is parsed into blocks, and blocks are compared per directory, so moving a
block between files or reformatting it is not counted as a change. A commit is \emph{usable} if it
changes between 2 and 30 blocks: one block gives no co-change to predict, and larger commits are bulk
refactorings whose co-changes are mostly coincidental (ROSE uses the same cut-off of
30~\cite{zimmermann2004rose}). The first 20\% of each repository's usable commits only feed the
history rules and produce no queries, so that history-based signals are not judged on an empty
history.

\paragraph{Dataset.} Table~\ref{tab:dataset} summarizes the 53 core repositories: 9,269 usable
commits (typically 5 changed blocks each) result in 55,541 queries with up to 200 existing candidates
each. On average, the seed is a new block in 38\% of a repository's queries, and 81\% of the true
answers survive the candidate pre-filter, which is the upper bound of the recall any method can
reach.

\begin{table}[ht]
\centering
\small
\caption{The 53 core repositories by kind.}
\label{tab:dataset}
\begin{tabular}{@{}lrrrr@{}}
\toprule
Kind & Repositories & Commits mined & Usable commits & Query commits \\
\midrule
Deployment      & 26 & 19,998 & 5,407 & 4,337 \\
Module library  & 21 & 11,332 & 2,827 & 2,268 \\
Educational     &  6 &  4,538 & 1,035 &   829 \\
\midrule
Total           & 53 & 35,868 & 9,269 & 7,434 \\
\bottomrule
\end{tabular}
\end{table}
